% =============================================================================
% cap2_hardware/02_requisitos_tipologias.tex
% Seção 2.2 -- Requisitos de aquisição (v2.2 — segunda compactação)
% =============================================================================
\section{Requisitos de aquisição}
\label{sec:requisitos-tipologias}

A definição do hardware de captura é tratada como problema de análise de viabilidade. As restrições mapeadas na Seção~\ref{sec:caracterizacao-pontos} são traduzidas em pré-requisitos que delimitam o espaço de soluções aceitáveis, separando os \textbf{mínimos} --- sem os quais o pipeline do Capítulo~\ref{cap:pipeline} é inviável --- dos \textbf{desejáveis}, que diferenciam uma solução adequada à escala e à rotina da AEX Gatosdoc2. A confrontação desses requisitos com as tipologias arquiteturais disponíveis é objeto da Seção~\ref{sec:tipologias-arquiteturais}.

Os \emph{requisitos mínimos} delimitam o que a infraestrutura precisa entregar para que o pipeline seja tecnicamente viável. O primeiro é o \textbf{acionamento autônomo por movimento}, realizável por duas vias funcionalmente análogas mas energeticamente distintas: por \emph{hardware}, com sensor infravermelho passivo (PIR) que mantém o subsistema principal em sono profundo até variação térmica no campo de visão; ou por \emph{software}, com câmera permanentemente ativa que dispara gravação ao detectar variação significativa de pixels entre quadros sucessivos. A literatura é convergente quanto à necessidade de algum mecanismo de disparo: \emph{ghost frames} acionados por vento, sombra ou variação térmica espúria respondem por metade ou mais do volume bruto capturado, e a gravação contínua é inviável em escala plurianual~\cite{bruce2025largescale,velez2022choosing,swanson2015snapshot}. A taxa de \emph{ghost frames} efetivamente entregue é parâmetro que o filtro de pré-processamento do pipeline precisará dimensionar.

Compõem ainda os mínimos a \textbf{vedação ambiental} (classe IP66 ou superior, Seção~\ref{sec:acondicionamento-geometria}), a \textbf{autonomia energética} compatível com pontos sem corrente alternada (Seção~\ref{sec:energia-conectividade}), o \textbf{armazenamento local} suficiente para o regime \emph{offline-first} entre ciclos de coleta e a \textbf{captura noturna por infravermelho próximo} (NIR), sem a qual mais de metade dos pontos perderia a janela crepuscular e noturna em que parte significativa da atividade da colônia ocorre. Cada um desses requisitos restringe não apenas o que o módulo precisa ser, mas o regime de qualidade de imagem com que o pipeline terá de operar.

Aos mínimos somam-se \emph{requisitos desejáveis}, que traduzem o diferencial de engenharia do projeto. O primeiro é a \textbf{capacidade de processamento na borda} (\emph{Edge AI}): inferências leves que descartam quadros sem animal, filtram vegetação em movimento e rejeitam fauna não-alvo antes da transmissão, desonerando o canal de comunicação, reduzindo armazenamento na nuvem e encurtando a malha de validação humana, conforme reportado em implantações recentes~\cite{velasco2024smartcameratraps,pishdast2025smartcameratraps,mulero2025addressing}. O segundo é a \textbf{conectividade nativa} para sincronização opcionalmente intermitente, que dispensa coleta manual sistemática de cartões e habilita o protocolo de validação humana no laço. Esses dois requisitos não condicionam a existência do sistema, mas determinam o quanto do pipeline pode ser deslocado para a borda em vez de operar inteiramente em servidor.
